%-------------------------
% Cover Letter — Akash Dubey -> Scale AI (AI Builder Intern)
% Generated from cover-letter-template.tex
%------------------------
\documentclass[letterpaper,11pt]{article}
\usepackage[empty]{fullpage}
\usepackage{titlesec}
\usepackage[usenames,dvipsnames]{color}
\usepackage{enumitem}
\usepackage[hidelinks]{hyperref}
\usepackage{fancyhdr}
\usepackage[english]{babel}
\usepackage{setspace}
\ifdefined\pdfgentounicode\input{glyphtounicode}\fi
\pagestyle{fancy}
\fancyhf{}
\renewcommand{\headrulewidth}{0pt}
\renewcommand{\footrulewidth}{0pt}
\addtolength{\oddsidemargin}{-0.5in}
\addtolength{\evensidemargin}{-0.5in}
\addtolength{\textwidth}{1in}
\addtolength{\topmargin}{-.5in}
\addtolength{\textheight}{1.0in}
\urlstyle{same}
\raggedright
\setlength{\parindent}{0pt}
\setlength{\parskip}{8pt}
\ifdefined\pdfgentounicode\pdfgentounicode=1\fi
%-------------------------------------------
\begin{document}

\begin{center}
    \textbf{\Huge \scshape Akash Dubey} \\ \vspace{2pt}
    \small (212)-814-3681 $|$ \href{mailto:akashdub06@gmail.com}{\underline{akashdub06@gmail.com}} $|$
    \href{https://akashdubey.me/}{\underline{akashdubey.me}} $|$
    \href{https://github.com/AkeBoss-tech}{\underline{github.com/AkeBoss-tech}} $|$
    \small U.S. Citizen
\end{center}
\vspace{4pt}
\hrule
\vspace{10pt}

June 22, 2026

\vspace{2pt}
Hiring Team, Scale AI --- Data \& Technology --- San Francisco, CA / New York, NY

\vspace{6pt}
Dear Hiring Team,

Scale AI is betting that the next wave of enterprise leverage will come from
agentic systems that do real work inside messy operational workflows, not from
AI demos that stop at the chat window. That is the exact problem I have been
building on, which is why I am applying for the AI Builder Intern role. The
posting's emphasis on shipping internal tools, automating real workflows, and
measuring value from day one feels much closer to how I already work than a
typical internship.

At New York Life, I built an AI text and voice agent for insurance claims work
used by clients and 12,000 customer service representatives, then helped ground
it by assembling a knowledge base from 13,000 documents across SharePoint,
Confluence, Jira, the public site, and the intranet. The hard part was not just
calling a model API; it was connecting fragmented systems, turning ambiguous
business logic into something an agent could act on, and making the output
useful enough for people to trust. In parallel, at Samaritan Scout, I built an
agentic extraction pipeline with OpenAI and Gemini structured outputs that
reduced manual data-entry work by more than 90\%, which reinforced the same
lesson: the best AI work is operational, measurable, and tightly tied to how a
team actually gets things done.

Beyond any single project, my real edge is how I work. I am fast and fluent
with modern AI tooling and use it to ship production systems rather than
presentations, and I pick up unfamiliar stacks and domains quickly. This year
alone I moved from enterprise insurance automation to RAG infrastructure to
CUDA robot motion planning without losing momentum. Scale's focus on building
internal AI products for real teams is exactly why this role stands out to me:
it offers the chance to work in an AI-native environment where judgment, speed,
and usefulness matter more than process theater.

I would love to help Scale AI ship agentic tools that save teams real time and
expand what they can do. Thank you for your time and consideration --- I would
welcome the chance to talk.

\vspace{6pt}
Sincerely, \\
\textbf{Akash Dubey}

\end{document}
